\documentclass[12pt]{article}
\usepackage{graphicx}
\usepackage{url}
\usepackage[dvipsnames]{xcolor}
\usepackage{enumitem}
\usepackage{caption}

% Page and typography setup
\usepackage[legalpaper, margin=2cm]{geometry}

% Math and units
\usepackage{amsmath,amssymb,mathtools}
\usepackage{siunitx}
\sisetup{separate-uncertainty=true}

% Figures, tables, and captions
\usepackage{graphicx}
\usepackage{booktabs}
\usepackage[font=small,labelfont=bf]{caption}

% Hyperlinks and clever references
\usepackage[hidelinks]{hyperref}
\usepackage{cleveref}
\usepackage{setspace}


% Title
\title{Classification and Segmentation of }
\author{
    \textbf{Mahdi Shakouri} \\[0.4em]
    University of British Columbia \\[0.6em]
    PHYS 310 - Machine Learning in Physics and Astronomy
}
\date{April 14th, 2026}

\begin{document}
\maketitle

\section*{Abstract}


\begin{figure}[!h]
    \centering
    % \includegraphics[width=0.6\textwidth]{C:/MohammadMahdi/University/Fifth Yr/PHYS 410/MP1/imgs/antennae.jpeg}
    \caption*{The Antennae Galaxies are a pair of interacting galaxies with
    dramatic tidal tails formed by their gravitational interaction. 
    Credit: NASA, ESA, and the Hubble Heritage Team (STScI/AURA).}
\end{figure}

\clearpage
\setcounter{tocdepth}{3}
\tableofcontents
\clearpage


\section{Introduction}


\section{Methods}



\subsection{Database Description}
The original DDSM (Digital Database for Screening Mammography) dataset on the Cancer Imaging Archives (CIA) created contains thousands of "normal" cases with no tumors and no masks. However, the CBIS-DDSM (Curated Breast Imaging Subset of DDSM) is a modernized, cleaned-up subset of that original database.

The dataset is organized into two main categories: "calc" for calcifications and "mass" for masses. 
Every single row in the mass and calc CSV files represents a confirmed abnormality, and every single one of those rows includes:


\subsection{Preprocessing}

\section{Discussion}


\section{Conclusion and Summary}

\section{Results}


\clearpage
% References section (thebibliography prints the "References" heading)
\begin{thebibliography}{9}

\bibitem{ref:example1}
Choptuik, M. (2025). Physics 210: Intro Computational Physics: Course Notes and Resources [Website].
Retrieved from \href{https://laplace.physics.ubc.ca/410/}{https://laplace.physics.ubc.ca/410/}

\bibitem{ref:example2}
Toomre, A., and Toomre, J. (1972). Galactic Bridges and Tails. The Astrophysical Journal, 178(3), 623666.

\end{thebibliography}

\end{document}